\section{HMT--MD Ontology of Particles}
\Needspace{7\baselineskip}
\subsection{Material genealogical state: extension, route, record, signature, and representation}
A complete HMT state is not a bare point. It is the tuple

\[
 \mathfrak x=
 (\Gamma,\mathcal L_\Gamma,\sigma_\Gamma,\eta_\Gamma,
 \mathcal R_\Gamma,\mathcal B_\Gamma,\mathcal C_\Gamma),
\]
where \(\Gamma\) is a surviving route, \(\mathcal L_\Gamma\) its genealogical record, \(\sigma_\Gamma\) its integer signature, \(\eta_\Gamma\) its harmonic refinement, \(\mathcal R_\Gamma\) its representation, \(\mathcal B_\Gamma\) its return memory, and \(\mathcal C_\Gamma\) its composition or boundary datum. Two states with the same visible word may be distinct if they differ in prefix, block, memory, carry, sheet, or Witt shadow.

\Needspace{7\baselineskip}
\subsection{Particle as a persistent local section of the dimensional genealogical state}
A particle is an equivalence class of persistent sections that satisfies four conditions:

\begin{enumerate}
\item compatible prolongation through the TPK cylinders;
\item phase return with state memory;
\item irreducible physical representation or minimal invariant block;
\item spectral stability under the coupling that observes it.
\end{enumerate}
A persistent section is a compatible family.

\[
 s=(s_K)_{K\ge K_0},\qquad
 p_{L,K}(s_L)=s_K\quad(L\ge K\ge K_0).
\]
Two sections \(s\) and \(s'\) representan the same particle when there exists a level \(K_1\) and a family compatible of simetrias internal \(g_K\in G_K\) tal that

\[
 s'_K=g_Ks_K\quad(K\ge K_1),
 \qquad
 p_{L,K}g_L=g_Kp_{L,K},
\]
and the observable characters of the genealogical record coincide. This relation identifies descriptions that differ by internal gauge, but it does not identify routes with distinct memory, carry, sheet, block, or holonomy. The particle is the class

\[
 [s]_{\rm pers}
 =\{s':s'\sim s\},
\]
not a single visible word nor a cell chosen after knowing its mass.

If \(\widetilde\Gamma_K\) is the extension at level \(K\), persistence is expressed by

\[
 p_{K+9,K}(\widetilde\Gamma_{K+9})
 =\widetilde\Gamma_K,\qquad
 g(K+9)=g(K),
\]
without requiring

\[
 \mathcal B_{K+9}=\mathcal B_K.
\]
The particle is not the returning phase, but the genealogical identity preserved through the change of state. Its mass measures the spectral rigidity of that persistence when it is coupled to a material realization.

\Needspace{7\baselineskip}
\subsection{Antiparticle as the conjugate image of a persistent section under sheet inversion}
The dodecaphonic conjugation is

\[
 C_M(\mathbf t)=-\operatorname{rev}(\mathbf t).
\]
The particle--antiparticle pair is formed by two orientations of the same genealogical band. If \(\chi\) denotes the sheet orientation,

\[
 (\mathbf t,\chi)\sim(C_M\mathbf t,-\chi).
\]
For a band energy

\[
 E_{\rm band}(\mathbf t,\chi)
 =E_+(\mathbf t)+\chi E_-(\mathbf t),
\]
with

\[
 E_\pm(\mathbf t)
 =\frac12\bigl(E_0(\mathbf t)\pm E_0(C_M\mathbf t)\bigr),
\]
we obtain

\[
 E_{\rm band}(C_M\mathbf t,-\chi)
 =E_{\rm band}(\mathbf t,\chi).
\]
The equality of masses is therefore a consequence of the operator's parity under conjugation, while charge and orientation change. Equality of widths additionally requires the decay operator to be conjugated by \(C_M\).

\Needspace{7\baselineskip}
\subsection{Mass as a spectral character of persistence of \(1\) a material route}
Mass is the spectral measure, on the line \(\mathcal L_M\), of the operator resulting from compressing the character, towers, bidirectional reader, and clock onto the persistent subspace selected by the coupling:

\[
 \widehat M_{P,a}
 =P_{P,a}\,
 \mathbf m_\star^{\rm ret}
 R_{\rm act}^{\widehat K_{P,a}/2}
 \widehat{\mathcal A}_{P,a}^{(0)}
 R_{\rm act}^{\widehat K_{P,a}/2}
 P_{P,a}.
\]
The general output is

\[
 \mu_{P,a}^{\rm mass}(B)
 =\operatorname{Tr}\!\left(
 \rho_{P,a}E_{\widehat M_{P,a}}(B)\right).
\]
A scalar mass \(m_P\) exists if and only if

\[
 P_{\operatorname{supp}\rho}\widehat M_{P,a}
 P_{\operatorname{supp}\rho}
 =m_P P_{\operatorname{supp}\rho},
\]
equivalently, if the variance of \(\widehat M_{P,a}\) in the prepared state is zero. This definition includes the rank-one electron, the multiplets forced by Schur's lemma, and isolated resonance poles, without conflating them.

\Needspace{7\baselineskip}
\subsection{Charge as a transported caliber character by the material section}
Charge is the oriented character of the odd layer of the genealogical record under conjugation. Let \(\sigma_{\rm odd}\) be the odd signature. Then

\[
 Q(C_M\Gamma)=-Q(\Gamma),\qquad
 Q(\Gamma)=\mathfrak q\!\left(\sigma_{\rm odd}(\Gamma)\right),
\]
where \(\mathfrak q\) is the homomorphism from the odd signature to the charge lattice of the realization. The even mass signature and the odd charge signature are distinct objects: a word may have zero linear sum and nevertheless carry a nonzero mass signature.

\Needspace{7\baselineskip}
\subsection{Spin as a representation of the rotational action on the material fiber}
Spin is the holonomy class with which a persistent section responds to a complete rotation of the orientation frame. The internal root

\[
 J^2=-I
\]
separate the two closures:

\[
 U(2\pi)=
 \begin{cases}
 +I,&\text{tensorial sector},\\
 -I,&\text{spinorial sector},
 \end{cases}
 \qquad
 U(4\pi)=I.
\]
An integer-spin boson lies in a representation that closes at \(2\pi\); a half-integer-spin fermion requires the double circuit \(4\pi\). The spin value is not deduced from an isolated cardinal count: it is the weight of the representation of the holonomy group on the persistent subspace. In particular:

\[
 s=0:\text{ scalar section},\qquad
 s=\tfrac12:\text{ spinorial band},
\]
\[
 s=1:\text{ transverse vector section},\qquad
 s=2:\text{ tensor component of a metric realization}.
\]
The last row defines a representation type; it does not require the existence of an elementary spin-two particle.

Chirality and helicity are subsequent and distinct data. An oriented involution \(\Gamma_{\rm ch}^2=I\) defines

\[
 P_L=\frac{I-\Gamma_{\rm ch}}2,
 \qquad
 P_R=\frac{I+\Gamma_{\rm ch}}2,
\]
whereas helicity is the spectral sign of the rotation generator along the direction of propagation. The former classifies sheets of the representation; the latter depends on the kinematic state. This separation allows a neutrino to be fermionic through its closure \(4\pi\), even though its weak coupling selects a single chirality.

\Needspace{7\baselineskip}
\subsection{Exchange symmetry induced quantum statistics}
Statistics is the rule of completion of the space of identical paths:

\[
 \mathcal F_-(\mathcal H)
 =\bigoplus_{n\ge0}\wedge^n\mathcal H,\qquad
 \mathcal F_+(\mathcal H)
 =\bigoplus_{n\ge0}\operatorname{Sym}^n\mathcal H.
\]
In the exterior completion, the CAR relations

\[
 \{a_i,a_j^\dagger\}=\delta_{ij},\qquad
 \{a_i,a_j\}=0
\]
imply

\[
 n_i^2=n_i
\]
and therefore Pauli exclusion and the Fermi--Dirac distribution. In the symmetric completion, the CCR relations permit multiple occupancy and lead to Bose--Einstein. Statistics are not inferred from mass: they follow from how the route is composed under exchange.

\Needspace{7\baselineskip}
\subsection{Color as residual representation of \(SU(3)\) on the quark routes}
Color is the residual triadic orbit of a quark section within the charge and flavor fiber. The thirty-six colored positions decompose as

\[
 12_R+12_G+12_B.
\]
The physical object is not one of those colors, but the triple and its action:

\[
 \mathcal H_q^{\rm color}
 =\mathcal H_R\oplus\mathcal H_G\oplus\mathcal H_B.
\]
If the mass operator commutes with the irreducible color action,

\[
 U_g\widehat M_qU_g^*=\widehat M_q,
\]
Schur's lemma forces a common mass in the triplet. Choosing a red, green, or blue cell as representative would destroy covariance.

\Needspace{7\baselineskip}
\subsection{Taste as a class of route within a material generation}
Flavor is an inequivalent coupling class within a single gauge representation, characterized by charge, generation, leaf, fine signature, memory, and interaction projector. Two flavors may share a coarse fiber and differ in the intertwiner coupling them to the reader.

Formally, if \(F\) is the common fiber, the flavors \(f_i\) are representations \(\pi_i\) for which

\[
 \operatorname{Hom}_{G_F}(V_{f_i},\ell^2(F))\ne0
\]
and whose projectors \(P_{f_i}\) are not conjugate under a gauge symmetry that remains physical. Flavour is not the depth of the atlas and does not reduce to a conventional name.

\Needspace{7\baselineskip}
\subsection{Material generation as a typed recurrent layer, stratified by return depth, fine signature, and reader spectrum}
A generation is a recurrent stratum of the same grammar of charge and representation, distinguished by return depth, fine signature, and reader spectrum. The generation label is physical only when there is a projector that intertwines that stratum with the coupling channel. A depth \(G_j\) must therefore not be identified automatically with the electron, muon, tau, or sterile neutrino.

\Needspace{7\baselineskip}
\subsection{Gauge field as a connection on the charge fibers}
A gauge field is the realization of a frame redundancy over the fibers. Its elementary bosons belong to the null direction of the mass operator while the symmetry remains exact. A gauge mass is not obtained by approximating a positive character to zero, but through a breaking, a longitudinal coupling, or a declared pole realization.

\Needspace{7\baselineskip}
\subsection{Composite particle as a persistent section of the product of linked paths}
A composite particle is a persistent section of the composition operator, not the sum of the masses or signatures of its constituents. If \(\Gamma_1,\ldots,\Gamma_n\) are constituent routes, first construct

\[
 \mathfrak C_{12}
 (\Gamma_1,\ldots,\Gamma_n;\mathcal T,\partial,\mathrm{carry})
 =\mathcal L_{\rm comp},
\]
where \(\mathcal T\) is the temporal tree and \(\partial\) the link boundary. Only afterwards does the reader act

\[
 \mathcal L_{\rm comp}\longrightarrow\sigma_{\rm comp}
 \longrightarrow\chi_{\rm comp}
 \longrightarrow\widehat M_{\rm comp}.
\]
The binding energy resides in the composition cocycle, boundary memory, and modification of the spectrum; it is not introduced as an external correction.

\Needspace{7\baselineskip}
\subsection{Resonance as a transient section and width as a decay rate}
A resonance is a persistent non-unitary mode of the coupled operator. If

\[
 \lambda=re^{-i\theta},\qquad 0<r<1,
\]
then

\[
 E=\frac{\hbar_{\rm int}}{t_0}
 \bigl(\theta-i\log r\bigr)
 =M-\frac{i}{2}\Gamma.
\]
The mass \(M\) and width \(\Gamma\) are the real and imaginary parts of a pole, but they require distinct readers. The width may also be obtained from a compressed semigroup:

\[
 T_X=Q_XUQ_X,\qquad
 \lambda_X=-\frac1{t_0}\log\rho(T_X),\qquad
 \Gamma_X=\hbar_{\rm int}\lambda_X.
\]
No width is deduced solely from the fact that a route possesses finite horizon.

\Needspace{7\baselineskip}
\subsection{Virtual section as an intermediate non-asymptotic extension of the atlas}
A section is virtual when its prolongation ends in a non-trivial manner:

\[
 \operatorname{Ext}_{L}(s)\ne\varnothing,\qquad
 \operatorname{Ext}_{L+1}(s)=\varnothing,\qquad
 \mathcal L(s)=L<\infty.
\]
Its primary content is the horizon \(L\), the terminal obstruction, the finite genealogical register, and the propagation amplitude. It is neither an imaginary-mass particle nor an as-yet-unclassified documentary state.

\Needspace{7\baselineskip}
